\uuid{3alB}
\titre{Logique}
\auteur{bodin}
\organisation{exo7}
\licence{CC BY-NC-SA 4.0}
\datecreate{2020-01-21}
\niveau{L1}
\module{Algèbre}
\chapitre{Logique, ensemble, raisonnement}
\sousChapitre{Autre}
\difficulte{}
\theme{Logique}
\exo7qcmid{arnaud.arnaud.010}

\contenu{
\begin{question}
\qclassification{100.99}
Soit $E$ un ensemble fini de cardinal $n$, soit $A\subset E$ un
ensemble à $p$ éléments, et  $B \subset E$ un ensemble à
$q$ éléments. On note $\mathcal{S} = \{ (a,b) \in A \times B \mid a\neq b \}$ et
 $\mathcal{T} = \{ (I,b) \text{ avec } I\subset A \mid \operatorname{Card} I =r \text{ et } b\in B \}$.
\begin{answers}
    \bad{Si $A \cap B = \emptyset$ alors $\operatorname{Card} \mathcal{S} = p+q$.}

    \good{Si $A \cap B = \emptyset$ alors $\operatorname{Card} \mathcal{S} = pq$.}

    \bad{Si $A \subset B$ alors $\mathcal{S} = \emptyset$.}

    \bad{$\operatorname{Card} \mathcal{T} = C_n^p \times r$.}

    \good{$\operatorname{Card} \mathcal{T} = C_p^r \times q$.}
\end{answers}
\end{question}
}
